% GENERATED FILE. Edit canonical lessons, then run npm run generate:books.
% canonical-source-hash: fnv1a64:1baf71698a957e73
% canonical-lessons: TA-C395-natattu, TA-C395-uruvakku, TA-C395-nikku, TA-C395-oli-abolish, TA-C395-kaivitu

\chapter{Taking Action}
\label{ch:ta-taking-action}
\hlchaptermodality{395}

\begin{chapteropening}
\textbf{By the end of this chapter:} \emph{I can talk about holding, creating, removing and giving up.}

Complete the last lesson of chapter 395: I can talk about holding, creating, removing and giving up.
\end{chapteropening}

\section[\texorpdfstring{\ta{நடத்து} (naṭattu)}{நடத்து (naṭattu)}]{\ta{நடத்து} (naṭattu) — to conduct; to hold; to run}

\label{lesson:TA-C395-natattu}



\begin{quote}
Before the new one: say the Tamil for on behalf of, then the Tamil for about.
\end{quote}

\subsection*{You'll want to know: \ta{நடத்து}}
\textbf{\ta{நடத்து}} — \emph{naṭattu} — \textquotedblleft{}to conduct; to hold; to run\textquotedblright{}.

\begin{quote}
\ta{நல்ல} \ta{வானிலை} \ta{இருந்தால்} \ta{பள்ளி} \ta{விளையாட்டுப்} \ta{போட்டி} \ta{நடத்தும்}.
\end{quote}

\begin{quote}
\emph{nalla vāṉilai iruntāl paḷḷi viḷaiyāṭṭup pōṭṭi naṭattum.}
\end{quote}

\begin{quote}
If the weather is good, the school will hold a sports meet.
\end{quote}

Past (I …): \ta{நடத்தினேன்}, \emph{naṭattiṉēṉ}.

\ta{கூட்டம்} \ta{நடத்து}, to hold a meeting. It is the causative of \ta{நட}, to walk: to make things walk.

\begin{cousinweb}
\textbf{\ta{நடத்து}} (\emph{naṭattu}) is native Dravidian: DEDR 3582 lists Tamil \emph{naṭa}, \textquotedblleft{}to walk, proceed, happen\textquotedblright{}, and its causative \emph{naṭattu}, \textquotedblleft{}to make walk, conduct, manage\textquotedblright{}.
\end{cousinweb}

\begin{grammarlens}[title={What you've built}]
One more word for saying what happens if.
\end{grammarlens}

\subsection*{Your turn}
\begin{itemize}
\raggedright
  \item \emph{Say it:} \emph{naṭattu}
  \item \emph{Say it:} \emph{naṭattu}, once more
  \item \emph{Say it:} \emph{nalla vāṉilai iruntāl paḷḷi viḷaiyāṭṭup pōṭṭi naṭattum.} — If the weather is good, the school will hold a sports meet.
  \item \emph{Recall:} say the Tamil for on behalf of, then the Tamil for about, then say \emph{naṭattu} again
\end{itemize}

\begin{tcolorbox}[breakable,colback=teal!4,colframe=teal!35!black,title={Before you move on}]
What does \ta{நடத்து} mean? (\textquotedblleft{}To conduct; to hold; to run\textquotedblright{}.) Say it once more.
\end{tcolorbox}
\section[\texorpdfstring{\ta{உருவாக்கு} (uruvākku)}{உருவாக்கு (uruvākku)}]{\ta{உருவாக்கு} (uruvākku) — to create; to form}

\label{lesson:TA-C395-uruvakku}



\begin{quote}
Before the new one: say the Tamil for about, then the Tamil for to conduct.
\end{quote}

\subsection*{You'll want to know: \ta{உருவாக்கு}}
\textbf{\ta{உருவாக்கு}} — \emph{uruvākku} — \textquotedblleft{}to create; to form\textquotedblright{}.

\begin{quote}
\ta{நாம்} \ta{ஒன்றாக} \ta{வேலை} \ta{செய்தால்} \ta{ஒரு} \ta{நல்ல} \ta{திட்டத்தை} \ta{உருவாக்கலாம்}.
\end{quote}

\begin{quote}
\emph{nām oṉṟāka vēlai ceytāl oru nalla tiṭṭattai uruvākkalām.}
\end{quote}

\begin{quote}
If we work together, we can create a good plan.
\end{quote}

Past (I …): \ta{உருவாக்கினேன்}, \emph{uruvākkiṉēṉ}.

\ta{உரு}, form + \ta{ஆக்கு}, to make. Its intransitive partner is \ta{உருவாகு}, to take shape.

\begin{cousinweb}
\textbf{\ta{உருவாக்கு}} (\emph{uruvākku}) is built from \emph{uru}, \textquotedblleft{}form\textquotedblright{}, borrowed from Sanskrit \emph{rūpa} (CDIAL 10803), and native \emph{ākku}, \textquotedblleft{}to make\textquotedblright{}, the causative of \emph{āku} (DEDR 333).
\end{cousinweb}

\begin{grammarlens}[title={What you've built}]
One more word for saying what happens if.
\end{grammarlens}

\subsection*{Your turn}
\begin{itemize}
\raggedright
  \item \emph{Say it:} \emph{uruvākku}
  \item \emph{Say it:} \emph{uruvākku}, once more
  \item \emph{Say it:} \emph{nām oṉṟāka vēlai ceytāl oru nalla tiṭṭattai uruvākkalām.} — If we work together, we can create a good plan.
  \item \emph{Recall:} say the Tamil for about, then the Tamil for to conduct, then say \emph{uruvākku} again
\end{itemize}

\begin{tcolorbox}[breakable,colback=teal!4,colframe=teal!35!black,title={Before you move on}]
What does \ta{உருவாக்கு} mean? (\textquotedblleft{}To create; to form\textquotedblright{}.) Say it once more.
\end{tcolorbox}
\section[\texorpdfstring{\ta{நீக்கு} (nīkku)}{நீக்கு (nīkku)}]{\ta{நீக்கு} (nīkku) — to remove; to dismiss}

\label{lesson:TA-C395-nikku}



\begin{quote}
Before the new one: say the Tamil for to conduct, then the Tamil for to create.
\end{quote}

\subsection*{You'll want to know: \ta{நீக்கு}}
\textbf{\ta{நீக்கு}} — \emph{nīkku} — \textquotedblleft{}to remove; to dismiss\textquotedblright{}.

\begin{quote}
\ta{தவறுகளை} \ta{நீக்கினால்} \ta{இந்தக்} \ta{கட்டுரை} \ta{நன்றாக} \ta{இருக்கும்}.
\end{quote}

\begin{quote}
\emph{tavaṟukaḷai nīkkiṉāl intak kaṭṭurai naṉṟāka irukkum.}
\end{quote}

\begin{quote}
If the mistakes are removed, this essay will be good.
\end{quote}

Past (I …): \ta{நீக்கினேன்}, \emph{nīkkiṉēṉ}.

What is removed takes -\ta{ஐ}; \ta{வேலையிலிருந்து} \ta{நீக்கு}, to dismiss from a job.

\begin{cousinweb}
\textbf{\ta{நீக்கு}} (\emph{nīkku}) is native Dravidian: DEDR 3685 lists Tamil \emph{nī}, \textquotedblleft{}to separate from, put away, abandon\textquotedblright{}, and \emph{nīkku}, \textquotedblleft{}to remove\textquotedblright{}.
\end{cousinweb}

\begin{grammarlens}[title={What you've built}]
One more word for saying what happens if.
\end{grammarlens}

\subsection*{Your turn}
\begin{itemize}
\raggedright
  \item \emph{Say it:} \emph{nīkku}
  \item \emph{Say it:} \emph{nīkku}, once more
  \item \emph{Say it:} \emph{tavaṟukaḷai nīkkiṉāl intak kaṭṭurai naṉṟāka irukkum.} — If the mistakes are removed, this essay will be good.
  \item \emph{Recall:} say the Tamil for to conduct, then the Tamil for to create, then say \emph{nīkku} again
\end{itemize}

\begin{tcolorbox}[breakable,colback=teal!4,colframe=teal!35!black,title={Before you move on}]
What does \ta{நீக்கு} mean? (\textquotedblleft{}To remove; to dismiss\textquotedblright{}.) Say it once more.
\end{tcolorbox}
\section[\texorpdfstring{\ta{ஒழி} (oḻi)}{ஒழி (oḻi)}]{\ta{ஒழி} (oḻi) — to abolish; to root out}

\label{lesson:TA-C395-oli-abolish}



\begin{quote}
Before the new one: say the Tamil for to create, then the Tamil for to remove.
\end{quote}

\subsection*{You'll want to know: \ta{ஒழி}}
\textbf{\ta{ஒழி}} — \emph{oḻi} — \textquotedblleft{}to abolish; to root out\textquotedblright{}.

\begin{quote}
\ta{எல்லோரும்} \ta{சேர்ந்தால்} \ta{ஊழலை} \ta{ஒழிக்கலாம்}.
\end{quote}

\begin{quote}
\emph{ellōrum cērntāl ūḻalai oḻikkalām.}
\end{quote}

\begin{quote}
If everyone joins together, we can root out corruption.
\end{quote}

Past (I …): \ta{ஒழித்தேன்}, \emph{oḻittēṉ}.

Transitive with -\ta{ஐ} and the past \ta{ஒழித்தேன்}; the intransitive \ta{ஒழி}, to cease, has \ta{ஒழிந்தது}.

\begin{cousinweb}
\textbf{\ta{ஒழி}} (\emph{oḻi}) is native Dravidian: DEDR 1009 lists Tamil \emph{oḻi}, \textquotedblleft{}to cease, be finished\textquotedblright{} with the past in -\ta{ந்த}-, and \textquotedblleft{}to finish, bring to an end\textquotedblright{} with the past in -\ta{த்த}-.
\end{cousinweb}

\begin{grammarlens}[title={What you've built}]
One more word for saying what happens if.
\end{grammarlens}

\subsection*{Your turn}
\begin{itemize}
\raggedright
  \item \emph{Say it:} \emph{oḻi}
  \item \emph{Say it:} \emph{oḻi}, once more
  \item \emph{Say it:} \emph{ellōrum cērntāl ūḻalai oḻikkalām.} — If everyone joins together, we can root out corruption.
  \item \emph{Recall:} say the Tamil for to create, then the Tamil for to remove, then say \emph{oḻi} again
\end{itemize}

\begin{tcolorbox}[breakable,colback=teal!4,colframe=teal!35!black,title={Before you move on}]
What does \ta{ஒழி} mean? (\textquotedblleft{}To abolish; to root out\textquotedblright{}.) Say it once more.
\end{tcolorbox}
\section[\texorpdfstring{\ta{கைவிடு} (kaiviṭu)}{கைவிடு (kaiviṭu)}]{\ta{கைவிடு} (kaiviṭu) — to abandon; to give up}

\label{lesson:TA-C395-kaivitu}



\begin{quote}
Before the new one: say the Tamil for to remove, then the Tamil for to abolish.
\end{quote}

\subsection*{You'll want to know: \ta{கைவிடு}}
\textbf{\ta{கைவிடு}} — \emph{kaiviṭu} — \textquotedblleft{}to abandon; to give up\textquotedblright{}.

\begin{quote}
\ta{மழை} \ta{பெய்ததால்} \ta{நாங்கள்} \ta{பயணத்தைக்} \ta{கைவிட்டோம்}.
\end{quote}

\begin{quote}
\emph{maḻai peytatāl nāṅkaḷ payaṇattaik kaiviṭṭōm.}
\end{quote}

\begin{quote}
Because it rained, we gave up the trip.
\end{quote}

Past (I …): \ta{கைவிட்டேன்}, \emph{kaiviṭṭēṉ}.

What is given up takes -\ta{ஐ}: \ta{பழக்கத்தைக்} \ta{கைவிடு}, to give up a habit.

\begin{cousinweb}
\textbf{\ta{கைவிடு}} (\emph{kaiviṭu}) is native: \emph{kai}, \textquotedblleft{}hand\textquotedblright{} (DEDR 2023), and \emph{viṭu}, \textquotedblleft{}to let go, leave\textquotedblright{} (DEDR 5393): to let something go from the hand.
\end{cousinweb}

\begin{grammarlens}[title={What you've built}]
One more word for saying what happens if.
\end{grammarlens}

\subsection*{Your turn}
\begin{itemize}
\raggedright
  \item \emph{Say it:} \emph{kaiviṭu}
  \item \emph{Say it:} \emph{kaiviṭu}, once more
  \item \emph{Say it:} \emph{maḻai peytatāl nāṅkaḷ payaṇattaik kaiviṭṭōm.} — Because it rained, we gave up the trip.
  \item \emph{Recall:} say the Tamil for to remove, then the Tamil for to abolish, then say \emph{kaiviṭu} again
\end{itemize}

\begin{tcolorbox}[breakable,colback=teal!4,colframe=teal!35!black,title={Before you move on}]
What does \ta{கைவிடு} mean? (\textquotedblleft{}To abandon; to give up\textquotedblright{}.) Say it once more.
\end{tcolorbox}
